
\section{1. Introduction}

\subsection{1.1 Motivation}

Full self-attention is $\mathcal{O}(L^2)$ in sequence length $L$, which is 
prohibitive for long-context inference. \textbf{Sliding-window attention} reduces 
this to $\mathcal{O}(L \cdot W)$ with $W \ll L$, but at a fundamental cost: 
tokens outside the window are invisible to each other. Information from the 
beginning of a long document cannot reach tokens near the end.

\textbf{Compressive memory} offers one solution. Instead of keeping every past 
token's key-value pair, maintain a small learned matrix $M \in \mathbb{R}^{d_c \times d_h}$ 
per head that summarizes the global context. Queries attend to this memory 
matrix as if it were a set of $d_c$ learned "virtual tokens".

\textbf{Infini-Attention} combines both ideas: local sliding-window attention for 
precise, position-sensitive computation, plus a learned compressive memory 
for global context. The two branches are combined through a learnable scalar 
gate:

\[
\text{output} = \text{LocalAttn}(Q, K, V) + \sigma(\beta) \cdot \text{MemAttn}(Q, M)
\tag{1.1}
\]

The gate $\sigma(\beta)$ is initialized close to zero, so the model starts by 
reproducing standard local attention and gradually learns to use the memory.

\subsection{1.2 Problem Statement}

Given a pretrained transformer $\mathcal{M}$ with $N_\ell$ attention layers, 
we seek a modified attention block $\mathcal{A}'$ such that:

1. $\mathcal{A}'$ reproduces the pretrained behavior at initialization.
2. $\mathcal{A}'$ reduces attention cost from $\mathcal{O}(L^2)$ to 
   $\mathcal{O}(L \cdot W + L \cdot d_c)$.
3. $\mathcal{A}'$ maintains access to global context via the memory matrix.
4. $\mathcal{A}'$ works across four attention architectures.

We do \textbf{not} aim to reduce model size. Infini-Attention targets 
\textbf{inference efficiency} for long sequences, not model compression.

\subsection{1.3 Our Approach: Infini-Attention}

Infini-Attention replaces the standard MHSA block with:

1. \textbf{Frozen base projections} $W_Q, W_K, W_V, W_O$ transferred from the 
   pretrained attention.
2. \textbf{Trainable LoRA adapters} on each projection.
3. \textbf{A learned compressive memory} $M$ per head.
4. \textbf{A learnable scalar gate} $\beta$.

The local branch uses sliding-window attention, which for causal models 
combines a causal mask and a distance-band mask:

\[
\text{LocalAttn}_{ij} = \text{softmax}\!\left(\frac{Q_i K_j^\top}{\sqrt{d_h}} + M_{ij}\right), \quad
M_{ij} = \begin{cases} 0 & j \leq i \text{ and } i - j < W \\ -\infty & \text{otherwise} \end{cases}
\tag{1.2}
\]

The memory branch computes:

\[
\text{MemAttn}_i = \text{softmax}\!\left(\frac{Q_i M^\top}{\sqrt{d_h}}\right) M
\tag{1.3}
\]

\subsection{1.4 Contributions}

\begin{itemize}
\item \textbf{A hybrid local+global attention block.} Infini-Attention combines
\end{itemize}
  sliding-window local attention with a learned compressive memory via a 
  gated scalar.
\begin{itemize}
\item \textbf{Multi-architecture support.} Same codebase handles flat, nested,
\end{itemize}
  causal, and GQA+RoPE models.
\begin{itemize}
\item \textbf{Initialization strategy.} The gate $\beta = -5$ ensures the memory
\end{itemize}
  starts nearly silent, preserving the pretrained local behavior.
\begin{itemize}
\item \textbf{Systematic evaluation.} Benchmarks on DistilBERT, RoBERTa, and GPT-2
\end{itemize}
  with ablations over $\beta$ initialization, compression dimension, and 
  window size.
\begin{itemize}
\item \textbf{Honest reporting.} Parameter count increases; the "compression" in the
\end{itemize}
  name refers to the memory, not the model.

\subsection{1.5 Paper Organization}

Section 2 reviews related work. Section 3 provides background. Section 4 
presents Infini-Attention. Section 5 details multi-architecture integration. 
Section 6 gives setup. Section 7 reports results. Section 8 analyzes 
contributions. Section 9 provides guidance. Section 10 lists limitations. 
Section 11 concludes.



\subsection{Notation}
